\scanpage{27}
C'est là une \uline{relation d'ordre} dans l'ens des Classes d'iso. de MD-graphes de type numérique $(g,\nu)$, ayant un plus grand élément. Dans la suite, nous allons préciser considérablement cette notion de spécialisation entre MD-graphes, en la reliant à une notion de \uline{morphismes} (de ``générisation'') de MD-graphes.

Notons que chaque $\mathcal M_{[G]}$ \deleted{\unsure{est donné}}\inserted{admet} un revêt\textsuperscript{t} fini $\mathcal M_G$, qui est représentable comme produit fini
\[
\mathcal M_G\simeq\prod_\alpha\underbrace{\mathcal M_{g_\alpha,\widehat I_\alpha}}_{\text{dim }3(g_\alpha-1)+\widehat\nu_\alpha}
\]
donc
\[
\deleted{\unsure{(26)}}(27)\quad\begin{aligned}
\dim\mathcal M_G/\mathbb Z&=\sum_{\alpha\in S}3(g_\alpha-1)+\sum\widehat\nu_\alpha\\
&=3\sum g_\alpha-3\underbrace{\operatorname{card}S}_{s_0}+\sum_{\alpha\in S}\underset{\substack{|\\\text{poids}\\\text{marqué}}}{\nu_\alpha}+\sum_{\alpha\in S}\underset{\substack{|\\\text{ordre}}}{\widetilde\mu_\alpha}\\
&=3\sum g_\alpha-3s_0+\nu+2s_1
\end{aligned}
\]
\mspage{26}
